% ============================================================
% Chapter 1 — Project Overview
% Source: PSYCON Engineering Design Document, Vol. 1, Ch. 1
% ============================================================
\section{Project Overview}
\label{sec:project-overview}

\subsection{Project Identification}
\label{sec:project-identification}

The system described in this document is named \textbf{PSYCON} (Psychological Condition Monitoring System). PSYCON is defined as a multimodal wearable system designed to continuously monitor physiological, environmental, and speech-derived indicators in order to assist in the early identification of changes associated with psychological well-being.

\subsection{Executive Summary}
\label{sec:executive-summary}

PSYCON is a two-module wearable research prototype that combines physiological sensing, environmental monitoring, and speech analysis into a unified platform for continuous mental health assessment. Rather than relying on a single biomarker or questionnaire, the system captures multiple complementary signals---including cardiovascular activity, electrodermal activity, body temperature, motion, ambient light, and speech---to provide a richer understanding of a user's psychological state. The platform is intended as an assistive screening and monitoring tool for research and educational purposes, and it is designed to support long-term data collection, personalized analysis, and future integration with AI-driven mental health technologies.

\subsection{Project Objective}
\label{sec:project-objective}

The objective of the project is to develop a wearable system capable of continuously collecting multimodal data that may help identify patterns associated with psychological conditions, while remaining portable, modular, low-cost, and suitable for research and science-competition demonstrations.

\subsection{Primary Goals}
\label{sec:primary-goals}

\begin{itemize}
    \item Develop a wearable prototype using commercially available components.
    \item Acquire synchronized physiological, environmental, and audio data.
    \item Support continuous monitoring over extended periods.
    \item Enable future AI-based analysis of multimodal signals.
    \item Maintain a modular architecture that allows independent development and testing of subsystems.
    \item Keep the design scalable for future PCB integration.
\end{itemize}

\subsection{Intended Applications}
\label{sec:intended-applications}

\begin{itemize}
    \item Research in multimodal mental health monitoring.
    \item Longitudinal wellness tracking.
    \item Educational demonstrations.
    \item AI dataset collection.
    \item Human--computer interaction research.
    \item Science fair and engineering competitions.
\end{itemize}

\subsection{High-Level Architecture}
\label{sec:high-level-architecture}

PSYCON consists of two independent wearable modules: a \textit{Wrist Module}, responsible for physiological data acquisition, and an \textit{Audio Module}, responsible for speech acquisition and environmental sensing. \Cref{fig:system-architecture} presents the overall system architecture, including sensor allocation, controller assignment, power subsystems, and the wireless data path to the host/analysis layer.

\begin{figure*}[t]
\centering
\begin{tikzpicture}[
    font=\small,
    module/.style={draw, rounded corners=3pt, fill=blue!6, minimum width=4.6cm, minimum height=1cm, align=center, line width=0.6pt},
    sensor/.style={draw, rounded corners=2pt, fill=gray!8, minimum width=3.6cm, minimum height=0.8cm, align=center, line width=0.5pt, font=\scriptsize},
    ctrl/.style={draw, rounded corners=2pt, fill=orange!12, minimum width=3.6cm, minimum height=0.8cm, align=center, line width=0.6pt, font=\scriptsize\bfseries},
    pwr/.style={draw, rounded corners=2pt, fill=green!10, minimum width=3.6cm, minimum height=0.8cm, align=center, line width=0.5pt, font=\scriptsize},
    host/.style={draw, rounded corners=3pt, fill=violet!8, minimum width=4.2cm, minimum height=1.2cm, align=center, line width=0.7pt},
    arrow/.style={-{Stealth[length=2.2mm]}, line width=0.6pt},
    dashedbox/.style={draw, dashed, rounded corners=4pt, inner sep=10pt, line width=0.5pt}
]

% ---- Wrist Module cluster ----
\node[sensor] (gsr) at (0,6.4) {GSR Electrodes\\(via ADS1115)};
\node[sensor] (max30102) at (0,5.4) {MAX30102\\PPG / Heart Rate};
\node[sensor] (mpu6050) at (0,4.4) {MPU6050\\IMU (Motion)};
\node[sensor] (mcp9808) at (0,3.4) {MCP9808\\Temperature};
\node[ctrl] (esp32wrist) at (0,2.2) {ESP32 DevKit\\(Wrist Controller)};
\node[pwr] (pwrwrist) at (0,1.0) {2$\times$500\,mAh LiPo (parallel)\\+ TP4056 USB-C Charger};

\node[dashedbox, fit=(gsr)(max30102)(mpu6050)(mcp9808)(esp32wrist)(pwrwrist), label={[font=\bfseries\small]above:Wrist Module}] (wristbox) {};

% ---- Audio Module cluster ----
\node[sensor] (mic) at (8.6,5.4) {INMP441\\I\textsuperscript{2}S MEMS Microphone};
\node[sensor] (light) at (8.6,4.4) {TEMT6000\\Ambient Light Sensor};
\node[ctrl] (esp32audio) at (8.6,3.2) {ESP32 DevKit\\(Audio Controller)};
\node[pwr] (pwraudio) at (8.6,2.0) {1$\times$500\,mAh LiPo\\+ TP4056 USB-C Charger};

\node[dashedbox, fit=(mic)(light)(esp32audio)(pwraudio), label={[font=\bfseries\small]above:Audio Module}] (audiobox) {};

% ---- Host / analysis layer ----
\node[host] (host) at (4.3,-0.6) {Host / Analysis Layer\\(Data Logging, Storage,\\Future AI-Based Analysis)};

% ---- Internal connections: Wrist Module ----
\draw[arrow] (gsr.south) -- (max30102.north);
\draw[arrow] (max30102.south) -- (mpu6050.north);
\draw[arrow] (mpu6050.south) -- (mcp9808.north);
\draw[arrow] (mcp9808.south) -- (esp32wrist.north);
\draw[arrow] (pwrwrist.north) -- (esp32wrist.south);

% ---- Internal connections: Audio Module ----
\draw[arrow] (mic.south) -- (esp32audio.north);
\draw[arrow] (light.south) -- (esp32audio.north);
\draw[arrow] (pwraudio.north) -- (esp32audio.south);

% ---- Wireless links to host ----
\draw[arrow, dashed] (esp32wrist.south) to[out=-70,in=160] node[pos=0.55, above, sloped, font=\scriptsize]{Wi-Fi / BLE} (host.west);
\draw[arrow, dashed] (esp32audio.south) to[out=-110,in=20] node[pos=0.55, above, sloped, font=\scriptsize]{Wi-Fi / BLE} (host.east);

\end{tikzpicture}
\caption{High-level system architecture of PSYCON, showing the Wrist Module (physiological sensing) and Audio Module (speech and environmental sensing) as independently powered ESP32-based subsystems that communicate wirelessly with a host/analysis layer.}
\label{fig:system-architecture}
\end{figure*}

\subsubsection{Wrist Module}
\label{sec:wrist-module}

\textbf{Purpose.} Physiological data acquisition.

\textbf{Sensors.}
\begin{itemize}
    \item GSR electrodes (via ADS1115)
    \item MAX30102 pulse oximeter / heart-rate sensor
    \item MPU6050 inertial measurement unit
    \item MCP9808 digital temperature sensor
\end{itemize}

\textbf{Controller.} ESP32 DevKit (30-pin).

\textbf{Power.} Two 500\,mAh LiPo batteries connected in parallel ($\approx$1000\,mAh total capacity), charged and protected via a TP4056 USB-C charging module.

\subsubsection{Audio Module}
\label{sec:audio-module}

\textbf{Purpose.} Speech acquisition and environmental sensing.

\textbf{Sensors.}
\begin{itemize}
    \item INMP441 I\textsuperscript{2}S MEMS microphone
    \item TEMT6000 ambient light sensor
\end{itemize}

\textbf{Controller.} ESP32 DevKit (30-pin).

\textbf{Power.} One 500\,mAh LiPo battery, charged and protected via a TP4056 USB-C charging module.

\subsection{Data Collected}
\label{sec:data-collected}

\Cref{tab:data-collected} summarizes the categories of data acquired by PSYCON across its two modules.

\begin{table}[H]
\centering
\caption{Data modalities collected by PSYCON.}
\label{tab:data-collected}
\small
\begin{tabularx}{\columnwidth}{@{}l Y@{}}
\toprule
\textbf{Category} & \textbf{Signals / Data} \\
\midrule
Physiological & Heart rate; pulse waveform (PPG); skin conductance (EDA/GSR); skin temperature; motion and orientation \\
Environmental & Ambient light level \\
Audio & Speech recordings for future AI feature extraction \\
\bottomrule
\end{tabularx}
\end{table}

\subsection{Design Philosophy}
\label{sec:design-philosophy}

The system follows five core engineering principles:

\begin{itemize}
    \item \textbf{Modularity} --- each subsystem can operate and be tested independently.
    \item \textbf{Scalability} --- prototype hardware should transition cleanly to a custom PCB.
    \item \textbf{Maintainability} --- use widely available, well-documented modules.
    \item \textbf{Low Cost} --- balance performance with affordability.
    \item \textbf{Research Focus} --- prioritize data quality and documentation over product miniaturization.
\end{itemize}

\subsection{Current Hardware Configuration}
\label{sec:current-hardware-configuration}

\Cref{tab:hardware-config} summarizes the current bill-of-materials-level hardware configuration for the PSYCON prototype, organized by controllers, power components, and sensors.

\begin{table}[H]
\centering
\caption{Current hardware configuration of the PSYCON prototype.}
\label{tab:hardware-config}
\small
\begin{tabularx}{\columnwidth}{@{}l Y c@{}}
\toprule
\textbf{Category} & \textbf{Component} & \textbf{Qty.} \\
\midrule
Controllers & ESP32 DevKit (30-pin) & 2 \\
\midrule
\multirow{2}{*}{Power} & TP4056 USB-C charger module (protected) & 2 \\
 & 500\,mAh 3.7\,V LiPo battery (HLCY 651735P) & 3 \\
\midrule
\multirow{6}{*}{Sensors} & MAX30102 pulse oximeter / heart-rate sensor & 1 \\
 & ADS1115 16-bit ADC & 1 \\
 & MPU6050 IMU & 1 \\
 & MCP9808 temperature sensor & 1 \\
 & TEMT6000 ambient light sensor & 1 \\
 & INMP441 I\textsuperscript{2}S MEMS microphone & 1 \\
 & GSR electrodes (analog front end TBD) & 1 set \\
\bottomrule
\end{tabularx}
\end{table}

\subsection{Communication Interfaces}
\label{sec:communication-interfaces}

\Cref{tab:comm-interfaces} lists the communication interfaces used to connect each sensor to its respective ESP32 controller.

\begin{table}[H]
\centering
\caption{Communication interfaces by sensor.}
\label{tab:comm-interfaces}
\small
\begin{tabularx}{\columnwidth}{@{}l Y@{}}
\toprule
\textbf{Interface} & \textbf{Connected Devices} \\
\midrule
I\textsuperscript{2}C & MAX30102, MPU6050, ADS1115, MCP9808 \\
I\textsuperscript{2}S & INMP441 \\
Analog & TEMT6000; GSR electrodes (through ADS1115) \\
Wi-Fi / BLE & ESP32 wireless communication (both modules) \\
\bottomrule
\end{tabularx}
\end{table}

\subsection{Engineering Status}
\label{sec:engineering-status}

\begin{validationbox}
\textbf{Confirmed:}
\begin{itemize}
    \item Component selection
    \item Module compatibility
    \item I\textsuperscript{2}C address compatibility
    \item Overall system architecture
    \item Power subsystem concept
    \item Sensor allocation between modules
\end{itemize}
\end{validationbox}

\begin{requirementbox}
\textbf{Pending validation:}
\begin{itemize}
    \item Final GSR analog front end
    \item Full firmware integration
    \item Electrical characterization (current, runtime, thermal)
    \item Long-duration testing
\end{itemize}
\end{requirementbox}

\subsection{Assumptions and Scope}
\label{sec:assumptions-scope}

This document distinguishes between three categories of technical claims:

\begin{itemize}
    \item \textbf{Confirmed information} --- derived from hardware, datasheets, and documented design decisions.
    \item \textbf{Engineering assumptions} --- based on standard module implementations.
    \item \textbf{Experimental items} --- items that require physical testing before they can be claimed as verified.
\end{itemize}

\begin{assumptionbox}
The distinction above is maintained throughout this document to preserve technical accuracy suitable for engineering documentation and science-competition review. Where a claim has not yet been experimentally validated, it is explicitly flagged rather than presented as a confirmed result.
\end{assumptionbox}
